% English counterpart of the focal Spanish insertion supplied by Ley 9 puertas.
\section{Horizon composition and thermal-geometric recovery}
\label{ins29:vii:horizonte}

The composition of the constants distinguishes what each physical
reading retains of the same generated state. We use the action and
length sections, the angular pair, and the constitutive velocity already
constructed from APP--TRIT--TPK. Abbreviate \(\eta=\eta_{\rm ret}\)
and denote by \(s_0\) the action section whose SI expression is
\(10^{-34}\,\mathrm{J\,s}\). Thus,
\[
 \hbar=\eta s_0,\qquad
 L_\alpha=r_N\alpha^{16}L_*,\quad r_N=\frac{5759}{23040},\qquad
 G=\frac{c^3L_\alpha^2}{\eta s_0}.
\]
The inscribed return determines \(r_N\); the local norm construction
determines the power \(16\); the section \(L_*\) retains the longitudinal
coordinatization. The function \(\eta=H_5-90\mathcal C_\pi/\pi\), with
its coefficients and regional dependencies, is the one obtained in
\eqref{md:compos:seccion-accion}. The following substitutions use these outputs
and retain the units and domains of their readers.

\subsection{Fixed area, fixed mass, and preservation of the thermal product}

For comparison with an Einstein-gravity horizon, the Bekenstein--Hawking
expression is \(S_{\rm BH}=k_Bc^3\mathcal A/(4G\hbar)\).
The coefficient \(1/4\) belongs here to this comparison law.
For a stationary Schwarzschild horizon without charge or rotation,
\(R_S=2GM/c^2\), \(\mathcal A=4\pi R_S^2\), and
\(k_BT_H=\hbar c/(4\pi R_S)\).

\begin{proposition}[Two constraints on horizon composition]
\label{ins29:vii:area_masa}
In the specified realization, the structural outputs give
\begin{align*}
 S_{\mathcal A}&=\frac{k_B\mathcal A}{4L_\alpha^2},&
 T_{\mathcal A}&=\frac{\eta s_0c}{2k_B\sqrt{\pi\mathcal A}},\\
 S_M&=\frac{4\pi k_Bc^2L_\alpha^2M^2}{\eta^2s_0^2},&
 T_M&=\frac{\eta^2s_0^2}{8\pi k_BL_\alpha^2M},
 \qquad S_MT_M=\frac{Mc^2}{2}.
\end{align*}
In a comparison between sections that retains \(c,L_\alpha,k_B\),
the fixed-area reading preserves \(S_{\mathcal A}\) and makes
\(T_{\mathcal A}\) proportional to \(\eta\); the fixed-mass reading
gives \(S_M\propto\eta^{-2}\) and \(T_M\propto\eta^2\).
\end{proposition}
\begin{proof}
The identity \(G\hbar=c^3L_\alpha^2\) immediately gives the first
entropy. Substituting \(R_S=\sqrt{\mathcal A/(4\pi)}\) into the
temperature gives the first temperature. At fixed mass,
\(R_S=2cL_\alpha^2M/(\eta s_0)\); its area and reciprocal produce
the other two formulas. Multiplication cancels the common factors,
leaving \(Mc^2/2\).
\end{proof}

The cancellation of action in the fixed-area entropy expresses a
compatibility between its gravitational and longitudinal readings.
At fixed mass, the same action section enters the horizon geometry,
and its dependence remains in the entropy. The balance must retain
the chosen constraint throughout the comparison.
The radial parameter \(R\) in the earlier realization
\(\hbar c/(2\pi R)\) equals \(2R_S\) in this comparison: both
expressions then represent the same temperature. This identification
retains the factor of two between geometric radius and acceleration scale.

\subsection{Angular separation and recovery of the action section}

The modular boundary realization retains eighteen tritic occupations
in each incidence sector. Write \(X=N_{90}\), \(Y=N_{120}\),
\(N=X+Y\), and \(M=X-Y\). Its specified separation is
\[
 \Delta_{\rm mod}=
 \frac{(A_{\deg}-C^*_{\deg})\pi/180}{9\cdot90}
 \left(1-\frac7{12}\frac\pi{729}\right),\qquad
 d=30\Delta_{\rm mod},\quad f_\pi=1-\frac{7\pi}{8748}.
\]
The coordinate \(d\) retains both the angular conversion \(\pi/180\)
and the incidence difference \(120-90=30\). In this realization, the
distribution over microscopic configurations is proportional to
\(\exp[-d(3X+4Y)]\).
We use \(A_{\deg}=1000\alpha\),
\(C^*_{\deg}=2(\eta+\alpha)\), and set
\(\xi=C^*_{\deg}/A_{\deg}=\tanh s\), on the positive-action branch
\(1/500<\xi<1\). The shape rapidity \(s\) is dimensionless; both
coordinates of the angular ratio are expressed in degrees.

\begin{proposition}[Compatibility curve and recovery reader]
\label{ins29:vii:curva}
With \(d_*=50\pi f_\pi\alpha/243\),
\[
 d=d_*(1-\xi)=d_*(1-\tanh s),\qquad
 0<d<\frac{499}{500}d_*.
\]
For \(m=1,\ldots,36\), the preceding distribution satisfies
\[
 \frac{\Pr(M=m)}{\Pr(M=-m)}=e^{dm}.
\]
Consequently, each probability ratio determines the same \(d\) and
recovers the coordinates
\[
 \xi=1-\frac d{d_*},\qquad
 \eta=499\alpha-\frac{2430d}{\pi f_\pi},\qquad
 C^*_{\deg}=1000\alpha\left(1-\frac d{d_*}\right).
\]
\end{proposition}
\begin{proof}
Substitution of the angular pair gives
\(d=\pi f_\pi(499\alpha-\eta)/2430\).
Since \(\eta=500\alpha\xi-\alpha\), this gives the first equality
and, from the specified branch, the interval. For the second identity,
each occupation configuration \((X,Y)\) has a counterpart \((Y,X)\)
with the same multiplicity. The ratio of their weights is
\(e^{-d(3X+4Y)+d(3Y+4X)}=e^{d(X-Y)}\).
Summing over \(X-Y=m\) retains that factor. Taking the logarithm and
solving for \(\xi\) and \(\eta\) gives the three recovery formulas.
\end{proof}

This composition establishes a joint check: the thirty-six logarithmic
ratios, divided by their respective \(m\), must give the same coordinate.
Recovery acts on already generated outputs and retains their causal order.
Its result connects the occupation bias to the elliptical shape and
action section, without fixing an additional radius or energy unit.

The channels of the Barbero functional are then recovered as
\[
 q_-=e^{-27d/f_\pi},\qquad
 q_+=D_A/q_-,\qquad D_A=e^{-100\pi\alpha/9}.
\]
Indeed, \(27d/f_\pi=\pi(A_{\deg}-C^*_{\deg})/180\), whereas
\(q_+q_-=D_A\). Substitution of these channels and \(C^*_{\deg}(d)\)
into \eqref{md:compos:funcional-barbero} determines \(\gamma(d)\).
The area coefficient then retains its composition
\(\gamma(d)a_0(d)L_\alpha^2\). For an area eigenvalue
\(\mathcal A_n=\gamma a_0L_\alpha^2n\), the horizon comparison gives
\(S_{\rm BH}(\mathcal A_n)/k_B=\gamma a_0n/4\).
The entropy of the reduced state additionally retains its probabilities
and bipartition. Both evaluations remain available with their respective
objects: horizon geometry and entanglement distribution.
